% =============================================================================
% CHAPTER 1: Introduction
% =============================================================================
% SPDX-FileCopyrightText: 2026 Frank Langbein <frank@langbein.org>
% SPDX-License-Identifier: LPPL-1.3c
%
% Suggested sections: Motivation and context; Aims and objectives; Contributions;
% Dissertation structure (roadmap). Optional: Scope (boundaries of the work),
% Terminology or key terms. Start with a short chapter intro (no \section) before
% the first section. This file also demonstrates template features; see
% Section~\ref{sec:ch1-reference}. Remove that section for your own report.
% =============================================================================

% -----------------------------------------------------------------------------
% Chapter intro (no section): one short paragraph before the first \section
% -----------------------------------------------------------------------------
% This chapter introduces the dissertation. It outlines the motivation and context for the work, the aims and objectives, the main contributions, and the structure of the remaining chapters. The following sections expand on each of these.

% -----------------------------------------------------------------------------
% Motivation and context (typical first section)
% -----------------------------------------------------------------------------
% Use \section, \subsection, \label{key}. Refer with \ref{key} or \autoref{key}.

\section{Motivation and Context}\label{sec:motivation}

This dissertation focusses on the game of Battleship Solitaire, not to be confused with the more commonly known two-player Battleship game. The solitaire version is a deterministic, logic-based puzzle game played by a single person. The aim is to reconstruct a solution on the grid using the given row/column tallies, and individual cell hints.

This is an NP-complete problem, which is the main motivation factor behind this project. As the grid size increases or the number of hints provided decreases, the search space explodes combinatorially.

Because it is NP-complete, brute-force guessing is mathematically infeasible for large boards. Therefore, the problem serves as ideal testbed for advanced optimisation techniques, such as \gls{ilp}, combined with modern-day constraint satisfaction solvers.

Existing literature proposes multiple ways to translate the rules of Solitaire Battleship into linear algebra. One method uses modelling based on the cells, one models based on the ships. The comparison of these sets the stage for this project. 

% \newthought{This template illustrates} the structure and style of a dissertation using the Qyber dissertation class. As in a real introduction, you would state the problem domain, why it matters, and how your work fits. Acronyms use the long-short style: first use gives the full form (e.g.\@ \gls{uk}); later use \acrshort{api} or \acrlong{api} when you need only short or long.

% As shown in previous work~\citep{example-paper}, we follow standard practice.
% %F Clarify whether this applies to all cases.
% \citet{example-book} provide a longer treatment; multiple citations: \citep{example-paper,example-conf}. You can link to other chapters (e.g.\@ Chapter~\ref{ch2}, Section~\ref{sec:objectives}) and to figures (Fig.~\ref{fig:demo}). The class provides \url{https://example.com} and \href{https://example.com}{link text} for URLs and hyperlinks.

% -----------------------------------------------------------------------------
% Aims and objectives (typical second section)
% -----------------------------------------------------------------------------
% Use \citep{} (parenthetical) or \citet{} (narrative), not \cite.

\section{Aims and Objectives}\label{sec:objectives}

The main aim of the project is to determine the most computationally efficient \gls{ilp} paradigm for solving the NP-Complete Battleship Solitaire problem through comparative analysis of two \gls{ilp} formulations, with one additional variant.

To achieve this aim successfully, the project will have the following objectives:

\begin{enumerate}[label=\textbf{O\arabic*:}, leftmargin=*]
  \item To define and implement two \gls{ilp} formulations for the Battleship ruleset. The first being a Cell-based (Big-M) constraint model, and the second a Ship-based (Set Packing) candidate model.
  \item To engineer a scalable, automated evaluation environment that integrates a puzzle generator, a standard \gls{csplib} format parser, and the Gurobi optimisation engine to facilitate standardised solver benchmarking.
  \item To construct a decoupled \gls{gui}, built on a layered architecture with thread-safe message passing between a Gurobi background solver and the CustomTkinter front-end, allowing for the real-time visualisation of solver outputs and the interactive construction of problems.
  \item To conduct a systematic evaluation comparing the computational overhead, search-tree pruning efficiency, and scalability limits of the formulated \gls{ilp} encodings against datasets of varying heuristic difficulty and grid dimensions. 
\end{enumerate}

% -----------------------------------------------------------------------------
% Contributions (typical third section)
% -----------------------------------------------------------------------------

\section{Contributions}\label{sec:contributions}

This dissertation presents a comprehensive algorithmic analysis of \gls{ilp} applied to Battleship Solitaire, an NP-Complete problem. The main contributions of this research are:

\begin{itemize}
  \item \textbf{Automated Evaluation Pipeline and Generator:} The development of a robust Python framework capable of parsing standard \gls{csplib} Prolog datasets, alongside a dynamic puzzle generator designed to create mathematically valid instances up to $30 \times 30$ in size for scalability analysis.
  \item \textbf{Comparative \gls{ilp} Formulations:} The formal mathematical specification and Gurobi implementation of two distinct solver paradigms: a Cell-based model utilizing Big-M boundary constraints, and a Ship-based model employing candidate generation and Edge-Based Set Packing.
  \item \textbf{Constraint Impact Analysis and Tuning:} A Constraint Impact Analysis revealed the 'Overhead vs. Pruning' trade-off, demonstrating that while explicit tally logic reduced solve times on hard instances by pruning dead-end search branches, redundant parity constraints introduced matrix bloat affected performance in certain conditions.
  \item \textbf{Scalability and Performance Discovery:} Conclusive evidence demonstrating that the Ship-based Set Packing formulation strictly dominates the Cell-based approach. The Ship model maintains an immense speed advantage, efficiently solving $30 \times 30$ zero-hint instances in under \SI{2}{\second}, where the Cell model severely weakens, taking an average of \textasciitilde\SI{11}{\second}.
\end{itemize}

% -----------------------------------------------------------------------------
% Dissertation structure / roadmap (typical final section of introduction)
% -----------------------------------------------------------------------------

\section{Dissertation structure}\label{sec:roadmap}

% Outline how the rest of the dissertation is organised so the reader knows what to expect. Reference chapters by label.

The remainder of this dissertation is organised as follows.

\textbf{Chapter~\ref{ch2}} establishes the background: the formal context of Battleship Solitaire as a discrete tomography problem, its complexity classification, and the prior \gls{ilp} formulations that motivate this work.

\textbf{Chapter~\ref{ch3}} presents the methodology, design and implementation. It formalises the puzzle, describes the layered software architecture, and gives the full mathematical specifications of the Cell-Based Model, the Ship-Based Model and the four valid inequalities of the Improved Cell Model. The chapter closes with the dataset construction strategy and the implementation details of the Gurobi mapping, the macOS thread-safety solution, and the headless evaluation pipeline.

\textbf{Chapter~\ref{ch4}} reports the evaluation. It defines the metrics, describes the three experimental scenarios, and presents the results across the standard $10 \times 10$ \gls{csplib} dataset and the \num{200}-instance generated dataset spanning grids up to $30 \times 30$. The discussion analyses dense-versus-sparse matrix structure, the role of presolve and root-node relaxation, and the scale-dependent ``Overhead vs. Pruning'' paradox. The chapter ends with a discussion of limitations and threats to validity.

\textbf{Chapter~\ref{ch5}} draws the conclusions, summarises the contributions, and outlines three avenues of future work: cross-solver benchmarking, a hybrid ILP formulation, and high-performance-computing stress testing.

\textbf{Chapter~\ref{ch6}} reflects on the learning gained from the project and identifies transferable habits for future practice.

The appendices contain supplementary material and additional results: Appendix~\ref{app:supp} provides the parser pattern, the Big-M derivation, and the candidate-generator output. Appendix~\ref{app:results} provides per-grid-size summary tables and the CSPLib coverage breakdown.

% Algorithms (algorithm, algpseudocode) appear in the List of Algorithms (Algorithm~\ref{alg:demo}). Use \texttt{[t]} or \texttt{[!t]} for top placement like figures/tables. For code snippets use verbatim; for wide tables or figures use \texttt{\string
% \begin{sidewaystable}} or \texttt{\string
%   \begin{sidewaysfigure}} (rotating).\footnote{Footnote example. Captions and footnotes are set at least 11\,pt by the class.} Placeholder prose in other chapters uses \texttt{\string\lipsum[n]}. Example: \lipsum[1]

%     Use \texttt{\string\afterpage\{\string\clearpage\}} (afterpage) when you need to clear after the current page.

%     \begin{algorithm}[t]
%       \caption{Example algorithm (algpseudocode). Alternative: algorithmic for traditional layout.}
%       \label{alg:demo}
%       \begin{algorithmic}[1]
%         \State \textbf{Input:} data $x$
%         \State \textbf{Output:} result $y$
%         \State $y \gets 0$
%         \For{$i = 1$ \textbf{to} $n$}
%         \State $y \gets y + x_i$
%         \EndFor
%         \State \Return $y$
%       \end{algorithmic}
%     \end{algorithm}

%     Short code (verbatim):
% \begin{verbatim}
%   x = 1
%   y = x + 1
% \end{verbatim}

    % -----------------------------------------------------------------------------
    % Template reference (for authors: where to find package demos and conventions)
    % -----------------------------------------------------------------------------

    % \section{Template reference}\label{sec:ch1-reference}

    % \paragraph{Font comparison.}
    % To compare the class font options (\texttt{sans}, \texttt{serif}, \texttt{hacker}), this paragraph uses the default body text, then explicit \textrm{serif}, \textsf{sans}, and \texttt{monospace}, with inline math $E = mc^2$ and a URL \url{https://example.org}. Rebuild with \texttt{font=sans}, \texttt{font=serif}, or \texttt{font=hacker} to see the chosen default; the other families show how headings, code, and captions look in each setup.

    % This introduction doubles as a template: the sections above (Motivation, Aims and objectives, Contributions, Dissertation structure) are the usual structure for a dissertation Chapter~1. Replace their content with your own. The same file also demonstrates:

    % \begin{itemize}
    %     \item \textbf{Citations and links:} \verb|\citep|/\verb|\citet| (natbib); link colour DarkBlue, cite/URL colour DarkGreen (class default).
    %   \item \textbf{\texttt{\string\newthought}:} Section-opening phrase in small caps with vertical space (e.g.\@ in Motivation). Letterspacing is applied by default.
    %     \item \textbf{Quote environments:} \verb|\begin{quotation}| \verb|...| \verb|\end{quotation}|, \verb|\begin{quote}| \verb|...| \verb|\end{quote}|, and \verb|\begin{verse}| \verb|...| \verb|\end{verse}| are restyled (smaller type, indents; verse centred) by the class.
    %     \item \textbf{Full-width content:} With class option \texttt{fullwidth}, use \verb|\begin{fullwidth}| \verb|...| \verb|\end{fullwidth}| for block content that extends into the margin, and \texttt{fullwidthfigure}/\texttt{fullwidthtable} for full-width floats (e.g.\@ Figure~\ref{fig:fullwidth}).
    %     \item \textbf{Epigraph:} \verb|\epigraph{quote}{source}| for chapter- or section-opening quotes (right-aligned, small type).
    %   \item \textbf{Other:} glossaries (long-short acronyms), amsmath/mathtools, textcomp, pifont, siunitx, relsize, mhchem, enumitem, graphicx/caption, float/stfloats, subcaption, tabularray/longtable, algorithm/algpseudocode, verbatim, afterpage, lipsum, rotating. Prefer parenthetical asides in the main text; use footnotes only sparingly (see \texttt{main.tex}).
    % \end{itemize}

    % Conventions for citations, spaces (\verb|~|, \verb|\@|), dashes, and reviewer comments (\verb|%ID |) are in the comment block at the top of \texttt{main.tex}. Class options (parskip/parindent, fullwidth, italic) are documented there too.
